%%%%%%%%%%%%%%%%%%%%%%%%%%%%%%%%%%%%%%%%%%%%%%%%%%%%%%%%%%%%%%%%%%%%%%%%%%%
\chapter{Query Instructions and Operational Guidance}
\label{ch:query-instructions}
%%%%%%%%%%%%%%%%%%%%%%%%%%%%%%%%%%%%%%%%%%%%%%%%%%%%%%%%%%%%%%%%%%%%%%%%%%%

%%%%%%%%%%%%%%%%%%%%%%%%%%%%%%%%%%%%%%%%%%%%%%%%%%%%%%%%%%%%%%%%%%%%%%%%%%%
\section{Why Query-Side Instructions Help}
%%%%%%%%%%%%%%%%%%%%%%%%%%%%%%%%%%%%%%%%%%%%%%%%%%%%%%%%%%%%%%%%%%%%%%%%%%%

The Qwen embedding cards include guidance such as:

\begin{quotation}
\textit{We recommend that developers customize the instruct according to their specific
scenarios, tasks, and languages. Our tests have shown that in most retrieval scenarios, not
using an instruct on the query side can lead to a drop in retrieval performance by
approximately 1\% to 5\%.}
\end{quotation}

Embedding models are not only learning generic ``semantic similarity.'' They are also learning
task-conditioned similarity. Those are different questions:

\begin{itemize}
    \item ``what texts are generally similar?''
    \item ``what documents answer this search query?''
    \item ``what code snippet is relevant to this bug report?''
    \item ``what multilingual passage best matches this English question?''
\end{itemize}

An instruction helps disambiguate the retrieval intent. Without an instruction, the model may
optimize for a vaguer notion of semantic closeness. With an instruction, the model can better
align the query embedding toward the retrieval objective you actually care about.

%%%%%%%%%%%%%%%%%%%%%%%%%%%%%%%%%%%%%%%%%%%%%%%%%%%%%%%%%%%%%%%%%%%%%%%%%%%
\section{The Basic Pattern}
%%%%%%%%%%%%%%%%%%%%%%%%%%%%%%%%%%%%%%%%%%%%%%%%%%%%%%%%%%%%%%%%%%%%%%%%%%%

\begin{itemize}
    \item Query side: add an instruction prefix
    \item Document side: embed the raw document or chunk (no instruction)
\end{itemize}

Conceptually:

\begin{Verbatim}
Query embedding input:
"Instruct: Given a web search query, retrieve relevant passages
that answer the query.
Query: how does KV cache reuse work in llama.cpp?"

Document embedding input:
"KV cache reuse in llama.cpp allows ..."
\end{Verbatim}

The instruction tells the model what ``relevant'' means for this query.

%%%%%%%%%%%%%%%%%%%%%%%%%%%%%%%%%%%%%%%%%%%%%%%%%%%%%%%%%%%%%%%%%%%%%%%%%%%
\section{Why Not Instruct Both Sides}
%%%%%%%%%%%%%%%%%%%%%%%%%%%%%%%%%%%%%%%%%%%%%%%%%%%%%%%%%%%%%%%%%%%%%%%%%%%

You usually do not want to add the same instruction to every document chunk because:

\begin{itemize}
    \item It adds noise and wasted tokens
    \item Document embeddings are often precomputed once and reused many times
    \item The task conditioning is usually needed most on the query side
\end{itemize}

The query is where the ambiguity lives.

%%%%%%%%%%%%%%%%%%%%%%%%%%%%%%%%%%%%%%%%%%%%%%%%%%%%%%%%%%%%%%%%%%%%%%%%%%%
\section{What ``Customize the Instruct'' Really Means}
%%%%%%%%%%%%%%%%%%%%%%%%%%%%%%%%%%%%%%%%%%%%%%%%%%%%%%%%%%%%%%%%%%%%%%%%%%%

It does not mean you need a different prompt for every single query. It usually means:

\begin{itemize}
    \item Choose a stable instruction template for each retrieval use case
    \item Adapt it to your language and domain
    \item Keep it consistent during indexing and evaluation
\end{itemize}

Examples of retrieval modes that may want different instructions:

\begin{itemize}
    \item Documentation search
    \item Code search
    \item Support-ticket retrieval
    \item Multilingual knowledge-base retrieval
    \item Legal or medical retrieval with specialized vocabulary
\end{itemize}

%%%%%%%%%%%%%%%%%%%%%%%%%%%%%%%%%%%%%%%%%%%%%%%%%%%%%%%%%%%%%%%%%%%%%%%%%%%
\section{Example Instruction Templates}
%%%%%%%%%%%%%%%%%%%%%%%%%%%%%%%%%%%%%%%%%%%%%%%%%%%%%%%%%%%%%%%%%%%%%%%%%%%

\textbf{General documentation retrieval:}

\begin{Verbatim}
Instruct: Given a question, retrieve passages that best answer it.
Query: {query}
\end{Verbatim}

\textbf{Code retrieval:}

\begin{Verbatim}
Instruct: Given a software engineering question, retrieve code or
documentation relevant to solving it.
Query: {query}
\end{Verbatim}

\textbf{Multilingual retrieval:}

\begin{Verbatim}
Instruct: Given a multilingual search query, retrieve passages with
the most relevant semantic meaning, even across languages.
Query: {query}
\end{Verbatim}

\textbf{Internal knowledge-base retrieval:}

\begin{Verbatim}
Instruct: Given an internal operations question, retrieve documents
that directly help answer the question.
Query: {query}
\end{Verbatim}

%%%%%%%%%%%%%%%%%%%%%%%%%%%%%%%%%%%%%%%%%%%%%%%%%%%%%%%%%%%%%%%%%%%%%%%%%%%
\section{What the Reported 1\%--5\% Drop Means}
%%%%%%%%%%%%%%%%%%%%%%%%%%%%%%%%%%%%%%%%%%%%%%%%%%%%%%%%%%%%%%%%%%%%%%%%%%%

That reported drop should be interpreted as: not catastrophic, but large enough to matter
in retrieval systems.

Why it matters: a few percentage points in retrieval can be the difference between the right
chunk appearing in the top~5 versus disappearing below the cutoff. Once a relevant document
misses the first-stage retrieval window, reranking cannot recover it. So even a modest drop
at the embedding stage can have a noticeable end-to-end effect on RAG quality.

%%%%%%%%%%%%%%%%%%%%%%%%%%%%%%%%%%%%%%%%%%%%%%%%%%%%%%%%%%%%%%%%%%%%%%%%%%%
\section{Operational Implication}
%%%%%%%%%%%%%%%%%%%%%%%%%%%%%%%%%%%%%%%%%%%%%%%%%%%%%%%%%%%%%%%%%%%%%%%%%%%

If you adopt a Qwen3 embedding model in this repo, you should plan for:

\begin{enumerate}
    \item A consistent query instruction template
    \item A retrieval-specific evaluation pass using your own corpus
    \item Measuring top-k recall before and after changing the instruction
\end{enumerate}

This matters more than chasing leaderboard decimals, because the best instruction for your
language mix, chunking strategy, domain vocabulary, and query style may differ from the
default examples in the model card.

\begin{notebox}[Practical rule for first deployment]
Start with a simple, stable query instruction. Do not instruct documents. Benchmark retrieval
quality on your own data. Only then tune the instruction wording. That is the low-complexity
path to capturing most of the benefit the Qwen guidance is pointing at.
\end{notebox}

%%%%%%%%%%%%%%%%%%%%%%%%%%%%%%%%%%%%%%%%%%%%%%%%%%%%%%%%%%%%%%%%%%%%%%%%%%%
\section{Operational Checklist}
%%%%%%%%%%%%%%%%%%%%%%%%%%%%%%%%%%%%%%%%%%%%%%%%%%%%%%%%%%%%%%%%%%%%%%%%%%%

\subsection{Before First Deployment}

\begin{itemize}
    \item Confirm the RTEB leaderboard entry for your chosen model (or confirm it is absent
          and accept the uncertainty)
    \item Run 50--100 representative query-document pairs through the model and check
          top-5 recall manually
    \item Choose a query instruction template appropriate to your domain and record it
    \item Decide on chunk size, overlap, and strategy before indexing (re-indexing later
          is expensive)
    \item Ensure the llama-server port for embeddings is distinct from the chat model port
\end{itemize}

\subsection{After Deployment}

\begin{itemize}
    \item Monitor embedding latency separately from generation latency
    \item Re-embed and re-index when you change the embedding model or chunk strategy
    \item When adding new document types (code, PDF, HTML), validate that the chunker handles
          them correctly before bulk ingestion
    \item If retrieval precision drops after a corpus update, check for stale vectors from
          deleted or updated documents
\end{itemize}

\subsection{Upgrading the Stack}

\begin{enumerate}
    \item \textbf{First step:} embed + basic vector search (current target)
    \item \textbf{Second step:} add a Qwen3 reranker for top-of-list quality
    \item \textbf{Third step:} add hybrid search (BM25 + vector with RRF) if exact-match
          queries are failing
    \item \textbf{Fourth step:} evaluate ColBERT or multi-vector retrieval for complex
          multi-aspect queries
\end{enumerate}

Each step should be gated on measurement: only add complexity when you have data showing it
helps. The first step---a correctly configured embedding server with a good query instruction
template---captures the majority of the improvement available.
